\documentclass[11pt,a4paper]{article}
\usepackage[margin=25mm]{geometry}
\usepackage[T1]{fontenc}
\usepackage{lmodern}
\usepackage{microtype}
\usepackage{booktabs}
\usepackage{longtable}
\usepackage{enumitem}
\usepackage{xcolor}
\usepackage{listings}
\usepackage{hyperref}
\usepackage{fancyhdr}
\usepackage{titlesec}

\definecolor{ink}{HTML}{18212A}
\definecolor{teal}{HTML}{147D70}
\definecolor{codebg}{HTML}{F3F5F7}
\definecolor{codegray}{HTML}{59636E}
\hypersetup{
  colorlinks=true,
  linkcolor=teal,
  urlcolor=teal,
  pdftitle={API Sentinel Project Guide},
  pdfauthor={API Sentinel Project}
}
\pagestyle{fancy}
\fancyhf{}
\lhead{API Sentinel}
\rhead{Project Guide}
\cfoot{\thepage}
\titleformat{\section}{\Large\bfseries\color{ink}}{\thesection}{0.7em}{}
\titleformat{\subsection}{\large\bfseries\color{teal}}{\thesubsection}{0.7em}{}
\lstset{
  basicstyle=\ttfamily\small,
  backgroundcolor=\color{codebg},
  commentstyle=\color{codegray},
  breaklines=true,
  frame=single,
  framerule=0pt,
  columns=fullflexible,
  keepspaces=true,
  showstringspaces=false,
  tabsize=4
}
\setlist[itemize]{leftmargin=1.5em}
\setlist[enumerate]{leftmargin=1.8em}

\title{\Huge\bfseries API Sentinel\\[0.4em]
\Large Product, Architecture, User Guide, and Security Review}
\author{Project documentation}
\date{November 2026}

\begin{document}
\maketitle
\tableofcontents
\newpage

\section{Executive summary}

\subsection{In everyday language}
API Sentinel watches the requests handled by an API that a team owns. A small
middleware component sends request metadata to a central dashboard. The platform
groups the observations into short time windows, highlights unusual behavior,
and can open an incident when error, latency, anomaly, or validated model-risk
signals meet its configured-in-code thresholds. Engineers can attach deployment
events, sanitized application logs, and recovery notes to help investigate.

It is useful when a team needs more context than a simple ``is the server up?''
check. It does not guarantee that it can predict every outage, identify a proven
root cause, or automatically repair production. Its current investigation report
is evidence-based and deterministic; it does not call an LLM.

\subsection{Technical summary}
The project is a multi-tenant web application built with FastAPI, PostgreSQL,
Redis, Dramatiq, React, and TypeScript. A Python ASGI middleware captures
metadata, batches it asynchronously, and authenticates to the collector with a
per-service key. PostgreSQL stores telemetry, per-minute feature windows,
incidents, model versions, logs, deployment events, runbooks, and notification
delivery state. Scheduled workers analyze organizations and dispatch alert
deliveries. Scikit-learn and XGBoost provide anomaly and predictive-model
capabilities; SHAP supplies model attributions. A LangGraph workflow assembles
an investigation report from retrieved runbooks and correlated evidence.

\begin{quote}
\textbf{Scope statement.} This guide documents what is implemented in the
repository as of November 2026 and distinguishes it from proposed future work.
The project is suitable for local development and a carefully configured
single-host staging deployment. It has not had a formal penetration test or
production-scale load certification.
\end{quote}

\section{Why a user may need it}

A basic uptime checker asks whether a health URL returned a response at one
moment. API Sentinel records request-level behavior, which can help answer
questions such as:
\begin{itemize}
  \item Did a service begin returning more server errors?
  \item Did its response latency or traffic pattern change?
  \item Which open incident is associated with this service?
  \item Was a deployment recorded shortly before the incident?
  \item Are there relevant recovery notes in the workspace runbooks?
\end{itemize}

The people most likely to use it are backend developers, SREs, DevOps engineers,
and engineering leads who need a shared view of service health and incident
evidence. It can also help a small team learn the mechanics of telemetry
collection, tenant isolation, anomaly detection, and incident workflows.

\subsection{Example applications}
\begin{description}
  \item[API reliability] Watch server-side latency, request volume, and 5xx
  responses for FastAPI services.
  \item[Release investigation] Compare an incident time with deployment events
  reported by a CI job.
  \item[Error context] Ingest selected, sanitized error messages and correlate
  them with an incident.
  \item[Learning and prototyping] Use the development traffic simulator to
  explore normal, degrading, cascading-error, and flash-crowd patterns.
  \item[Workspace operations] Give a small engineering team a tenant-scoped
  dashboard, incident list, runbooks, and optional Slack notifications.
\end{description}

\section{Feature status: implemented and future work}

\begin{longtable}{p{0.25\linewidth}p{0.23\linewidth}p{0.44\linewidth}}
\toprule
\textbf{Area} & \textbf{Status} & \textbf{Current behavior} \\
\midrule
Authentication & Implemented & Email/password registration and sign-in;
HTTP-only session cookie; workspace owner role. \\
Tenant separation & Implemented & Organization filters and PostgreSQL row-level
security on tenant data tables; cross-organization behavior was exercised in
an isolated end-to-end run. \\
Telemetry & Implemented & Python FastAPI middleware and service-key ingestion
endpoints; request/response bodies are not captured. \\
Feature analysis & Implemented & Per-minute feature windows, anomaly scoring,
manual analysis endpoint, and optional minute scheduler. \\
Failure prediction & Implemented with a quality gate & Chronological holdout;
models remain candidates until the configured validation checks pass. \\
Incidents & Implemented & Threshold, anomaly, or validated-risk signals can
create an incident; users can change incident status. \\
Evidence investigation & Implemented, bounded & Telemetry, ingested logs,
deployment events, and keyword-matched runbooks are assembled by a LangGraph
workflow. Correlation is a clue, not causal proof. \\
Alerts & Implemented, configurable & Slack incoming webhook storage, test
delivery, queued incident delivery, and bounded retries. A real workspace
webhook must be supplied by the operator. \\
External provider sync & Not implemented & CI/CD and hosted log-provider data
must be posted to the ingestion endpoints by the user's integration. \\
LLM and vector search & Not implemented & Runbook retrieval is keyword-based;
there is no connected LLM or pgvector store in the current implementation. \\
Automated remediation & Not implemented & Investigation reports require human
review; no production action is executed. \\
Production hosting & Prepared, not deployed & A single-host Compose deployment
template exists. Domain, TLS termination, host, secrets, and operational
ownership must be configured by the operator. \\
\bottomrule
\end{longtable}

\section{System architecture}

\subsection{Runtime component view}
\begin{lstlisting}
Monitored FastAPI service
  request -> SentinelMiddleware -> application response
                  |
                  +-- bounded in-memory queue, batched HTTP POST
                           |
                           v
                     FastAPI collector <------ browser dashboard
                           |                         |
                           v                         v
                    PostgreSQL <-------------- React + TypeScript
                           ^
                           |
             Dramatiq worker <--- Redis <--- scheduler
                 |       |
                 |       +-- Slack webhook delivery
                 +-- feature analysis, training, incidents

Incident investigation reads telemetry, ingested logs, deployments,
and workspace runbooks, then returns a human-reviewed evidence report.
\end{lstlisting}

\subsection{Responsibilities by layer}
\begin{description}
  \item[SDK / collection] Measures route template, method, status, duration,
  content lengths, validated request/trace identifiers, and exception type. It
  avoids reading request and response bodies and authorization headers.
  \item[HTTP API] Handles account and workspace operations, service keys,
  telemetry, logs, deployments, incidents, runbooks, intelligence, and alerts.
  \item[PostgreSQL] Persists tenant data and serves as the durable source of
  truth. Telemetry event identifiers and integration event identifiers are
  deduplicated.
  \item[Redis / Dramatiq] Holds the background queue and ephemeral scheduling
  locks and heartbeat values. Redis is not the durable record of incidents.
  \item[Scheduler] Enqueues analysis and pending alert work every minute and
  model evaluation daily when enabled.
  \item[Worker] Processes queued jobs with retry limits and records a heartbeat.
  \item[Frontend] React and TypeScript UI. Vite proxies API calls in local
  development; production Nginx serves the SPA and forwards same-origin API
  paths to FastAPI.
\end{description}

\subsection{Data ownership and isolation}
Each service belongs to one organization. Its ingestion key is shown to the
owner once and stored in PostgreSQL as a SHA-256 hash. Collector calls identify
the service from the key; the server derives the organization from that service
record rather than trusting an organization identifier in the event payload.
User-facing endpoints derive the active organization from the signed-in
membership. Tenant-owned tables use PostgreSQL row-level security policies in
addition to application-level organization filters.

The production Compose template creates a runtime database role with
\texttt{NOSUPERUSER} and \texttt{NOBYPASSRLS}. The local development Compose
database uses its development owner role and should not be exposed as a
production database.

\section{End-to-end workflow}

\begin{enumerate}
  \item An owner creates an account and workspace.
  \item The owner creates a service and stores the one-time ingestion key in the
  monitored API's secret environment.
  \item The middleware observes HTTP metadata and puts events in its bounded
  asynchronous queue.
  \item The SDK sends batches to \texttt{/api/telemetry/batch} with the
  \texttt{X-Sentinel-Key} header. The collector validates the key, deduplicates
  events, and writes them to PostgreSQL.
  \item The scheduler queues feature analysis. A user can also run analysis
  from Intelligence.
  \item Per-minute observations become feature windows. Anomaly scoring uses
  Isolation Forest at 30 or more windows and a robust rolling baseline for
  shorter histories.
  \item Current error, latency, anomaly, or validated model-risk signals can
  create an incident. Default implemented thresholds include a 5\% 5xx rate,
  p95 latency of 1500 ms, anomaly score 0.65, and validated risk score 0.70.
  These values are implementation defaults, not a user-facing policy editor.
  \item The investigation workflow can correlate deployment events from the
  prior 24 hours, application logs around the incident, and matching workspace
  runbooks. Reports label associations as clues and require human review.
  \item If a Slack destination is configured, a durable delivery row is queued
  and the background worker delivers a concise incident message with bounded
  retries.
\end{enumerate}

\subsection{Model workflow}
The failure label asks whether the following ten minutes contain at least a
5\% 5xx rate or p95 latency of at least 1500 ms. Training uses time-ordered data
with a temporal purge between training and holdout windows. The classifier is
kept as a candidate unless enough positive and negative holdout examples exist
and it beats the configured baseline checks. In the isolated verification run,
training returned a candidate with 1,439 labeled windows and
\texttt{needs\_more\_validation\_data}; it was not activated. This is a safe
model outcome, not a system error.

\section{Tutorial: use the local application}

\subsection{Start the services}
From the project root, run PostgreSQL and Redis. Start the backend API, worker,
scheduler, and frontend in separate terminals following these commands:

\begin{lstlisting}[language=bash]
docker compose up -d postgres redis

cd backend
.venv/bin/alembic upgrade head
.venv/bin/uvicorn app.main:app --reload
# In another backend terminal:
.venv/bin/dramatiq app.tasks --processes 1 --threads 2
# In another backend terminal:
.venv/bin/python -m app.scheduler

# In another terminal from the repository root:
cd frontend
npm run dev
\end{lstlisting}

Open \url{http://localhost:5173}. Register or sign in. In Settings, confirm
Redis, scheduler, and worker are online.

\subsection{Create and instrument a service}
Create a service from the Services page and copy its ingestion key. Store the key
as \texttt{API\_SENTINEL\_KEY} in the monitored API's environment, then install
the SDK from the repository:

\begin{lstlisting}[language=bash]
pip install -e /path/to/api-sentinel/sdk-python
export API_SENTINEL_KEY='paste-the-service-key-in-your-terminal-env'
\end{lstlisting}

Register the middleware in the monitored FastAPI application's own entry point.
Do not paste the application setup into the SDK's \texttt{middleware.py} file.

\begin{lstlisting}[language=Python]
import os
from fastapi import FastAPI
from api_sentinel import SentinelMiddleware

app = FastAPI()
app.add_middleware(
    SentinelMiddleware,
    service_key=os.environ["API_SENTINEL_KEY"],
    collector_url="http://127.0.0.1:8000/api/telemetry/batch",
)
\end{lstlisting}

Restart the monitored application and send it requests. For example, API
Sentinel's local backend is instrumented for self-monitoring when its environment
contains a valid service key:

\begin{lstlisting}[language=bash]
curl -i http://127.0.0.1:8000/api/health
\end{lstlisting}

The middleware transmits metadata in the background. A successful collector
response is HTTP 202. Collector paths are excluded from instrumentation to avoid
a telemetry feedback loop. Normal health checks validate ingestion but do not
usually cause incidents because they return HTTP 200.

\subsection{Analyze, investigate, and alert}
\begin{enumerate}
  \item Open the Dashboard and select a time range.
  \item Open Intelligence and choose \textit{Analyze last 24 hours}.
  \item Review feature windows, model status, and any created incidents.
  \item Open an incident and run the investigation. Treat deployment-time
  correlation as context, not proof of cause.
  \item Add recovery instructions under Runbooks.
  \item As a workspace owner, connect a Slack incoming webhook in Settings and
  use the test-delivery action. Use only a webhook from the intended test
  workspace when validating.
\end{enumerate}

\subsection{Synthetic traffic}
The simulator is for a separate development service. It permanently writes
events to the selected service and can affect labels and incident history.
Cascade traffic spread over a 24-hour history may leave the final minute too
sparse to create an immediate incident. To exercise incident creation reliably,
send a concentrated burst of failing requests to a disposable service, then run
analysis. Do not use generated failure traffic in a production workspace.

\section{Evidence and API integration points}

\begin{longtable}{p{0.34\linewidth}p{0.56\linewidth}}
\toprule
\textbf{Route} & \textbf{Purpose} \\
\midrule
\texttt{POST /api/telemetry} & Ingest one request event using a service key. \\
\texttt{POST /api/telemetry/batch} & Ingest a batch of up to 100 events. \\
\texttt{POST /api/deployments} & Ingest one deployment event for incident correlation. \\
\texttt{POST /api/logs/batch} & Ingest up to 100 selected error/warning log records;
common credential patterns are redacted before persistence. \\
\texttt{POST /api/intelligence/analyze} & Analyze recent telemetry for the signed-in workspace. \\
\texttt{POST /api/intelligence/train} & Train or evaluate the workspace failure model. \\
\texttt{GET /api/intelligence} & Read model and feature-window summary. \\
\texttt{POST /api/incidents/\{id\}/investigate} & Create an evidence report for an incident. \\
\texttt{GET /api/operations/status} & Read Redis, scheduler, and worker status after sign-in. \\
\bottomrule
\end{longtable}

All evidence integrations use the service's \texttt{X-Sentinel-Key}. The product
does not pull data from GitHub, a CI provider, or a hosted log platform by itself;
those systems must post events to these routes.

\section{Deployment and operations}

\subsection{Prepared single-host deployment}
The repository includes \texttt{docker-compose.production.yml}, backend and
frontend Dockerfiles, an Nginx SPA/API proxy, a one-shot restricted-role setup
service, and a runtime-role migration. The web port is the only published application port;
PostgreSQL and Redis have no host port mapping. A TLS reverse proxy must sit in
front of Nginx because the included Nginx configuration serves HTTP only.

Copy \texttt{.env.production.example} to the ignored
\texttt{.env.production}, generate distinct 32-byte random values for the
administrator password, runtime database password, and session/encryption key,
then run:

\begin{lstlisting}[language=bash]
docker compose --env-file .env.production \\
  -f docker-compose.production.yml up -d --build
docker compose --env-file .env.production \\
  -f docker-compose.production.yml ps
\end{lstlisting}

The \texttt{sentinel\_app} role is created as a non-superuser without
\texttt{BYPASSRLS}. Migrations run separately with the administrator role; API
and worker processes use the restricted role. This template has been syntax
checked, built, and exercised in an isolated local stack. It still must be
deployed and checked on the selected host before public launch.

\subsection{Backups and recovery}
Run \texttt{backend/deploy/backup-production.sh} with a protected off-host
destination. It creates a PostgreSQL custom-format dump and sets file permissions
to owner-only. It does not encrypt the file. Encrypt backups at rest, set a
retention policy, and rehearse a restore into a disposable database. Keep
PostgreSQL and Redis volumes persistent; Redis queue state can be rebuilt only
for pending work, while PostgreSQL contains the durable business data.

\subsection{Operational checks}
\begin{itemize}
  \item Check \texttt{/api/health} and container health checks.
  \item Check Settings for Redis, scheduler, and worker heartbeats.
  \item Review API, worker, scheduler, and reverse-proxy logs.
  \item Alert on failed health checks, repeated notification failures, disk
  pressure, and missed backups using the hosting provider's monitoring.
  \item Restore a backup on a regular schedule and record the recovery time.
\end{itemize}

The repository does not include an external error-monitoring account, automatic
backup scheduler, or managed hosting configuration. Those depend on the chosen
provider and operator credentials.

\section{Security and vulnerability review}

\subsection{Implemented safeguards}
\begin{itemize}
  \item Ingestion keys are generated with high entropy and only their SHA-256
  hashes are stored in the database.
  \item Browser sessions use HTTP-only cookies, a 12-hour expiration, and
  SameSite=Lax. Production Compose sets secure cookies; local development does
  not.
  \item Tenant data queries use organization filters and row-level security on
  tenant tables. The production deployment uses a role that cannot bypass RLS.
  \item The middleware records metadata and does not inspect request/response
  bodies or authorization headers.
  \item Log messages are length-limited and common bearer, password, token,
  secret, and API-key patterns are redacted. This is not a complete sensitive
  data detector.
  \item Slack webhook URLs are restricted to the Slack webhook host and
  encrypted at rest using the configured application secret.
  \item The production Nginx template limits request body size and sets common
  browser security headers.
\end{itemize}

\subsection{Risks and recommended actions}
\begin{longtable}{p{0.18\linewidth}p{0.29\linewidth}p{0.43\linewidth}}
\toprule
\textbf{Priority} & \textbf{Risk} & \textbf{Action before public launch} \\
\midrule
High & A service key was previously pasted into source code. The SDK file was
cleaned and the current key is held in the ignored, mode-600 environment file. &
Rotate any key that appeared in source control, logs, tickets, or shared chat;
verify the old key is rejected and the replacement works. Never commit
\texttt{.env} files. \\
High & Public registration has no rate limit or email verification; login
failures are throttled per email in Redis but fail open if Redis is unavailable. &
Add signup abuse controls, email verification, and monitoring before allowing
unrestricted public account creation. \\
High & Production is not yet deployed behind verified TLS and a managed host. &
Use HTTPS end-to-end at the edge, set \texttt{COOKIE\_SECURE=true}, restrict
ingress, and configure trusted proxy headers. \\
Medium & Session tokens are stateless. Clearing the cookie on logout does not
revoke a copied token before its 12-hour expiry. & Add server-side session
revocation or shorten the accepted lifetime and document incident response. \\
Medium & Log redaction uses a small set of regular expressions; arbitrary PII or
secrets may remain in messages. & Sanitize at the source, minimize log fields,
and review retention and access before enabling log ingestion. \\
Medium & The normal Compose template has no automated backup encryption,
retention, or restore rehearsal. & Encrypt off-host backups, restrict access, and
prove restore procedures before launch. \\
Medium & Login failures are limited per email, but registration and general
ingestion have no distributed rate limit. The 1 MB payload cap is at the Nginx
edge. & Do not publish the API container directly; add signup abuse controls and
edge rate limits before internet exposure. \\
Medium & The app role still reads user and membership tables needed for login
and workspace selection; not every table has RLS. & Keep queries parameterized,
maintain explicit tenant filters, and add a focused review of identity-table
policies before multi-tenant scale-up. \\
Low & Model scores and investigation correlations can be wrong or incomplete. &
Keep human review, monitor false positives/negatives, and avoid describing
correlation as causal proof. \\
Low & npm reported no known vulnerabilities during the image build, but a
Python software-composition scan and formal penetration test have not been
completed. Python transitive packages are not hash-locked. & Run an authenticated
software-composition scan, lock transitive dependencies, and arrange an
independent security assessment before handling sensitive production telemetry. \\
\bottomrule
\end{longtable}

This is a source-based engineering review, not a penetration test or a guarantee
that the system is free of vulnerabilities. In this work session the CodeRabbit
CLI reported an authentication/Git-provider access error after its review
request was submitted, so no CodeRabbit findings were available. Local review,
the isolated E2E run, Python compilation, frontend production build, migration,
and Compose configuration checks were used instead.

\subsection{Threat model at a glance}
\begin{description}
  \item[Untrusted telemetry sender] May try a forged or stolen key, duplicate
  events, or submit oversized data. The collector authenticates a service key,
  bounds batch length and fields, and deduplicates event IDs; edge request
  limits are still important.
  \item[Malicious tenant] May try to access another workspace's records.
  Organization filters and RLS are tested; production must use the restricted
  role rather than a superuser.
  \item[Compromised account] May change workspace settings, create services, or
  configure outgoing Slack alerts according to role. There is no MFA or email
  verification in this version.
  \item[Compromised application secret] Can permit forged sessions and prevent
  decryption of stored Slack destinations if rotated without a re-encryption
  plan. Back up and rotate it deliberately.
  \item[Untrusted log content] May contain credentials or personal information.
  The server redacts common forms, but the source application remains responsible
  for data minimization.
\end{description}

\section{End-to-end verification record}

The recent isolated verification used a separate temporary PostgreSQL database,
a separate API port, and a separate Redis database index. It exercised:
\begin{itemize}
  \item anonymous request rejection, registration, sign-in, and service creation;
  \item key-required telemetry, valid batch ingestion, and duplicate handling;
  \item synthetic traffic, feature analysis, and deterministic incident creation;
  \item deployment and log ingestion, investigation evidence, and secret redaction;
  \item two organizations seeing only their own services and cross-organization
  incident mutation returning 404;
  \item candidate model evaluation and scheduler/worker/Redis heartbeats.
\end{itemize}

The candidate model remained inactive because the quality gate required stronger
holdout validation. The test database and its known schedule-lock keys were
removed afterward; the normal API and frontend remained healthy. Slack delivery
to a real workspace and browser-driven UI automation were not part of that run.

\section{Glossary}
\begin{description}
  \item[API] A program interface accessed over HTTP.
  \item[Telemetry] Request metadata used to describe service behavior.
  \item[Ingestion key] A secret assigned to one service so it can report events.
  \item[Feature window] A set of measurements aggregated for a short time span.
  \item[Anomaly] A behavior pattern that differs from the measured baseline.
  \item[Candidate model] A trained model that has not passed the activation gate.
  \item[Row-level security] PostgreSQL rules that restrict which rows a database
  role can read or modify in a query.
  \item[Runbook] Human-written operational steps for diagnosing or recovering
  from a known problem.
  \item[MTTR] Mean time to recovery, a common measure of how quickly service is
  restored after an incident.
\end{description}

\section{Project file map}
\begin{longtable}{p{0.42\linewidth}p{0.5\linewidth}}
\toprule
\textbf{Path} & \textbf{Role} \\
\midrule
\texttt{backend/app/main.py} & HTTP API, authentication, ingestion, and routes. \\
\texttt{backend/app/operations.py} & Feature analysis, incident creation, and model operations. \\
\texttt{backend/app/intelligence.py} & Feature engineering, anomaly scoring, and model utilities. \\
\texttt{backend/app/notifications.py} & Slack message delivery and retry state processing. \\
\texttt{backend/app/investigation.py} & LangGraph evidence-report workflow. \\
\texttt{backend/app/scheduler.py} & Minute scheduler and Redis schedule locks. \\
\texttt{backend/app/tasks.py} & Dramatiq background actors. \\
\texttt{sdk-python/api\_sentinel/middleware.py} & Monitored-service ASGI middleware. \\
\texttt{frontend/src/App.tsx} & React pages and user interactions. \\
\texttt{docker-compose.production.yml} & Single-host deployment template. \\
\texttt{backend/deploy/backup-production.sh} & PostgreSQL backup helper. \\
\texttt{backend/migrations/versions} & Database schema history. \\
\bottomrule
\end{longtable}

\section{Build this document}
The source is \texttt{PROJECT\_GUIDE.tex}. With a LaTeX distribution installed,
compile twice to populate the table of contents:

\begin{lstlisting}[language=bash]
pdflatex PROJECT_GUIDE.tex
pdflatex PROJECT_GUIDE.tex
\end{lstlisting}

The compilation toolchain is not bundled with the application repository.

\end{document}
